%! Author = Sujal Singh
%! Date = 10/23/23

% Preamble
\documentclass[11pt,a4paper,oneside]{scrarticle}

% Packages
\usepackage{mathptmx}
\usepackage{tabularray}
\usepackage{graphicx}
\usepackage{sectsty}
\usepackage{titlesec}
\usepackage{setspace}
\usepackage{romannum}
\usepackage{hyperref}

\renewcommand*{\maketitle}{
    \pagenumbering{gobble}
    \begin{center}
        \vspace*{8em}
        \includegraphics[width=200pt]{../../tools/pdf/template/logo} \\[40pt]
        \textbf{\Huge Law Assignment} \\[40pt]
        \begin{table}[htb]
            \label{tab:name-slip}
            \begin{tblr}{rows={30pt},colspec={|X|X[3]|},hlines}
                \textbf{Course}            & Industrial Internet of Things \\
                \textbf{Batch}             & 1                             \\
                \textbf{Name}              & Sujal Singh                   \\
                \textbf{Enrollment Number} & XXXXXXXXXXXX                  \\
            \end{tblr}
        \end{table}
    \end{center}
    \newpage
    \pagenumbering{arabic}
}

\linespread{1.5}
\titleformat{\section}{\Large\bfseries\centering}{\thesection.}{0.5em}{}

% Document
\begin{document}
    \maketitle

    \begin{center}
        \textbf{\LARGE The Indian Government's Stance on Fundamental Rights in the Digital World.}
        \\[10pt]
    \end{center}
    \begin{flushright}
        \textit{--- Sujal Singh} \\[20pt]
    \end{flushright}

    %------------------------------------------------------------------------------------------------------------------%
    %------------------------------------------------------------------------------------------------------------------%

    \section*{Abstract}\label{sec:abstract}
    The Internet is one of mankind's greatest achievements and an open Internet, free of censorship and control by any
    authority, is crucial to avoid a dystopian future, but ever since it's inception, government authorities have
    employed various tactics to censor the information published via the World Wide Web.\ This report details the
    state of fundamental rights as they exist in the digital landscape of India and how the Indian government has often
    violated these rights to censor the Internet and use it as a tool for mass-surveillance.\ Such acts go against
    constitutional values as envisioned by the authors of the Indian Constitution, and people in relevant fields of
    studies should feel responsible to oppose the continued decline of the open Internet, and they should push for
    technical solutions that uphold constitutional values.
\end{document}